\documentclass[11pt,a4paper]{article}
\usepackage{gtu-el-yazisi}
\begin{document}
\notbilgi{Makine Öğrenmesi}{Naive Bayes, kernel density ve k-nn}{Ders-5.pdf}

\sayfa{1}
\begin{tabular}{@{}l@{\quad}l@{}}
\textcolor{kirmizikalem}{\underline{Spam mail}} & \textcolor{kirmizikalem}{\underline{Naive Bayes}} \\[3pt]
kazandınız & $\rightarrow$ \textcolor{turuncukalem}{spam} \\
indirim & $\rightarrow$ \textcolor{turuncukalem}{spam} \\
proje & $\rightarrow$ \textcolor{turuncukalem}{ham} \\
merhaba & $\rightarrow$ \textcolor{turuncukalem}{ham}
\end{tabular}

kelimeler birbirine bağlı değilse
\[ P(x_1|x_2) = \frac{P(x_1,x_2)}{P(x_2)} = \frac{P(x_1)P(x_2)}{P(x_2)} = P(x_1) \]

\[ P(\text{spam}|x_1,\ldots,x_N) \overset{?}{\geq} P(\text{ham}|x_1,\ldots,x_N) \]
\[ P(y|x_1,\ldots,x_N) = \frac{\overbrace{P(x_1,\ldots,x_N|y)}^{\text{likelihood}}\;\overbrace{P(y)}^{\text{prior distrubution}}}{P(x_1,\ldots,x_N)} \]

\begin{align*}
\underbrace{P(x_1,\ldots,x_N,y)}_{\text{joint distribution}} &= P(x_1|x_2,\ldots,x_N,y)\,P(x_2,\ldots,x_N,y) \\
&= P(x_1|x_2,\ldots,x_N,y)\,P(x_2|x_3,\ldots,x_N,y)\,P(x_3,\ldots,x_N,y) \\
&= P(x_1|x_2,\ldots,x_N,y)\,P(x_2|x_3,\ldots,x_N,y)\cdots P(x_N|y)
\end{align*}
kelimeler bağımsızsa $\Rightarrow\;=P(x_1|y)\,P(x_2|y)\cdots P(x_N|y)$

\[ \frac{P(x_1,\ldots,x_N|\text{spam})\,{\color{kirmizikalem}P(\text{spam})}}{P(x_1,\ldots,x_N|\text{ham})\,{\color{kirmizikalem}P(\text{ham})}} \geq 1 \]
$\uparrow$ Bu \okunamadi{önsel} olacak

\textcolor{kirmizikalem}{$\ast$ $P(C_1)$}

$\ast$ $P(x_1|C_i) = \dfrac{P(x_1,C_i)}{P(C_i)}$

Burda bernolli dağılım kullandık, gauss de kullanılabilirdi

scikit learn naive bayes sınıflandırıcı

\sayfa{2}
\subsection{Kitapta non parametrik kısımlara gireceğiz}
\elcizimi[0.65]{sekiller/Ders-5-s2-1.png}[2]
olasılık dağılımı için kullanılabilir o zaman $\int f(x) = 1$ olmalıdır.

\elcizimi[0.5]{sekiller/Ders-5-s2-2.png}[2]
\[ \frac{1}{C}\sum_{i=1}^{N} K\!\left(\frac{x-x_i}{h}\right) = f(x) \]
$\rightarrow$ varyansa göre \qquad varyans büyükse $h$ \okunamadi{küçülür}

kernel Approximation \qquad fonksiyonlar birbirinin aynısı \\
window functions

\elcizimi[0.6]{sekiller/Ders-5-s2-3.png}[2]

Burda bütün verileri tutmamız lazım k-nn de \\
distance metriği ile çözersek
\[ d = |x - x_i| \]

\elcizimi[0.4]{sekiller/Ders-5-s2-4.png}[2]
kernel fonksiyonu \quad $k = \dfrac{1}{|x-x_i|}$

\okunamadi{1nn $k^{nn}$}
\[ \left\{ \max_i \frac{1}{|x-x_i|} \right. \]
2-nn için en büyük ikincisi

\sayfa{3}
\[ \min_i (k_i) \]

make blobs

Kernel density

Pesudocode soruları

\end{document}
